\texttt{std::allocator} "--- это аллокатор по умолчанию для всех контейнеров STL, который по сути является абстракцией. Он использует глобальные операторы \texttt{new} и \texttt{delete} для выделения и освобождения памяти, которые в свою очередь обращаются к функциям выделения памяти операционной системы. \texttt{std::allocator} "--- шаблонный класс с методами:

\begin{itemize}
	\item create() "--- создаёт аллокатор и отдаёт ему в распоряжение некоторый объём памяти;
	\item allocate() "--- выделяет память;
	\item deallocate() "--- освобождает память;
	\item construct() "--- создаёт объект в выделенной памяти;
	\item destroy() — уничтожает объект.
\end{itemize}

При выполнения запроса аллокации выполняются следующие шаги.

\begin{enumerate}
	\item Проверка размера: запрашиваемый объём памяти не должен превышать максимально возможный размер для объектов данного типа.
	\item Вызов глобальной функции \texttt{::operator new}.
	\item Системный запрос: обращение к куче операционной системы. Как правило, здесь выполняется вызов \texttt{malloc}, однако могут применяться и системные механизмы, такие как mmap (для Linux) \cite{itmo-assignment}. 
	\item Возврат указателя: операционная система возвращает указатель на выделенный блок памяти и передаёт его обратно аллокатору.
	\item Размещение объекта: вызов функции construct(), которая вызывает конструктор объекта по указанному адресу.
\end{enumerate}

Ниже приведён пример использования стандартного аллокатора C++.

\small
\begin{minted}{cpp}
#include <memory>
#include <vector>

using namespace std;

int main() {
    allocator<int> alloc;
	// Выделение памяти под 5 объектов типа int (создание массива)
    int* arr = alloc.allocate(5);
    // Создание объектов в выделенной памяти
    for (int i = 0; i < 5; i++) {
        alloc.construct(arr + i, i);
    }
    // Уничтожение объектов
    for (int i = 0; i < 5; i++) {
        alloc.destroy(arr + i);
    }
    // Освобождение памяти
    alloc.deallocate(arr, 5);
	// Создание вектора, который при операциях с ним (push_back, pop_back)
	// будет использовать данный аллокатор
    vector<int, allocator<int>> v;
}
\end{minted}
\normalsize

При решении большинства классических задач разработки нас устраивает стандартный аллокатор C++, и потребность в разработке собственного аллокатора возникает нечасто. Тем не менее, ручное управление памятью, постоянные вызовы \texttt{new} и \texttt{delete} могут привести к уязвимостям, сбоям и падению производительности. Основные проблемы, связанные с аллокацией памяти, включают в себя:

\begin{enumerate}
	\item Фрагментацию памяти. Куча разбивается на мелкие свободные блоки, из"=за этого при запросе на выделение большого непрерывного фрагмента памяти происходит сбой выделения, даже если общий объем свободной памяти достаточен.
	\item Снижение производительности. Динамическая аллокация требует времени, ситуация усугубляется в многопоточных программах, где требуется синхронизация между потоками для предотвращения гонки данных (data race).
	\item Утечки памяти. В C++ отсутствует автоматический сборщик мусора, поэтому если выделенная память не освобождается разработчиком своевременно, это приводит к исчерпанию системных ресурсов и, возможно, даже зависанию операционной системы. 
	\item Висячие указатели (dangling pointers, доступ к памяти, которая уже была освобождена), двойное освобождение (double free, попытка освободить один и тот же блок памяти дважды). В отличие от других проблем, которые лишь замедляют работу программы, эти ошибки могут привести к непредсказуемому поведению или аварийному завершению работы.
\end{enumerate}

Стандартный аллокатор при каждом вызове \texttt{new} или \texttt{delete} обращается к операционной системе, в отличие от пользовательских аллокаторов, работающих внутри уже выделенной памяти процесса. В аллокаторах, которые будут рассмотрены в дальнейшем, предварительно одним запросом к операционной системе выделяется большой непрерывный блок памяти, а затем уже из него с помощью средств языка программирования пользователю отдаются более мелкие блоки памяти под конкретные объекты. Это позволяет ускорить работу с памятью и сократить её фрагментацию \cite{pre-allocation}. В этой работе будут рассмотрены несколько базовых типов аллокаторов: пулловый, линейный и стековый.
